\documentclass[12pt]{article}
\usepackage[utf8]{inputenc}
\usepackage[T1]{fontenc}
\usepackage{hyperref}
\title{Etulo Language Resources for K-12 Curriculum Development}
\author{Compiled by Claude Sonnet 5 (claude-sonnet-5) for Andr\'as Kornai}
\date{}
\begin{document}
\maketitle

\section{Introduction}
Etulo (ISO 639-3: utr) is a Niger-Congo (Idomoid) language spoken by a minority community in Benue and Taraba states, Nigeria. Searches found one substantial scholarly resource, a 2025 grammar, and only second-hand mentions of other materials. No primers, textbooks, dictionaries or school programmes were found.

\section{Primary Educational Resources}
No Etulo primers or textbooks were found in search results. The grammar's bibliography reportedly lists unpublished 1975 notes, \emph{Let us Read in Etulo}, from the Nigerian Bible Translation Centre's Etulo Development Committee. I could not confirm that they are accessible. In the interview below, the grammar's author says learners cannot yet readily access structured online or physical learning resources.

\section{Linguistic and Reference Materials}
\begin{itemize}
\item Chikelu Ezenwafor-Afuecheta, \emph{Grammar of Etulo: A Niger-Congo (Idomoid) Language} (Open Book Publishers, 2025), a free open-access grammar. The author presents it as a foundation for pedagogical materials. \url{https://openbookcollective.org/books/book/5825837}
\item Catalogue record for the same grammar. \url{https://thoth-arch.lib.cam.ac.uk/handle/1811/921}
\item Interview with the grammar's author (Philological Society blog). \url{https://blog.philsoc.org.uk/2025/11/18/interview-with-chikelu-i-ezenwafor-afuecheta-the-author-of-etulo-a-niger-congo-idomoid-language/}
\end{itemize}
No dictionary was found.

\section{Religious and Community Resources}
According to the interview above, the Nigerian Bible Translation Trust, working with Etulo speakers, has developed an orthography proposal and translated parts of the New Testament. No details or access point were found; the Trust would need to be contacted directly.

\section{Digital and Audio Resources}
No audio, video, corpus or software resources for Etulo were found.

\section{Curriculum Development Context}
The 2025 grammar and the proposed orthography are the best starting points for materials development. Primers, readers and a dictionary would have to be created from scratch. Collaboration with Etulo speakers and the Nigerian Bible Translation Trust is the obvious route.
\end{document}